% AI-written, AI cross-reviewed; no human sign-off; not a claim that RH is proved
% Atlas unit P032. Local manuscript v3, item 7.9, heading on printed PDF p. 77.
% UpperTheory.V.ZK.sClassOf
% paper2/inputs_one.tex:377–432. Verbatim excerpt; SHA-256 270207101f563eb23aa0c1798c475898de913c6db3d696cbf6edc756a7b0ef0e.
% Original TeX uses the macros and cross-reference keys of the source manuscript.

\begin{formalstatement}[An integral cycle with uniform residue evaluations]\label{pn:P032}
Assume functoriality and the boundary identities for chains and prisms,
the prism composition law, the specified basic bar cycle and its
coefficient relation, the interpolation distribution, and the Steinitz
transport hypothesis. In any layer as above take \(E=e+1\), \(e\ge0\),
\(m=2E+1<\ell^c\), and \(a=\kappa^E\). For any \(\alpha\in F^\times\)
satisfying the difference condition there are \(N\), a degree-\(m\)
bar cycle \(y\) on \(\mathrm{GL}_N(\mathcal O_L)\) over
\(\mathbb Z/q\), and \(0\ne b\in\mathcal O_F\) such that the following
holds. For every reduction \(r:\mathcal O_F\to k\) to a finite field
which is onto on \(\mathcal O_L\) and has \(r(b)\ne0\), for every
character \(v:k^\times\to\mathbb Z/q\), and for every additive
evaluating family \(\theta\) satisfying the scalar law for \(r(\zeta)\),
there is a residue \(u\in k^\times\) of \(N_a(\alpha)\) for which
\[
 \theta\bigl(\operatorname{red}_r y\bigr)=
 \prod_{\substack{0\le j<m\\j\ne E}}(a-\kappa^j)\,v(u)
 \quad\text{in }\mathbb Z/q.
\]
\end{formalstatement}\nodestatus{P032}
\begin{outline}
We recall the construction in Blueprint Z2.K2~\bploc{sec:P-r4-K2}. Write
\(p(Y)=\prod_{0\le j<m,\,j\ne E}(Y-\kappa^j)=\sum_i B_iY^i\)
and form \(\beta=\prod_i(\gamma^{-i}\alpha)^{B_i}\).
Writing \((Y-a)p(Y)=\sum A_iY^i\), substitution of the difference
condition gives \(\operatorname{ord}_P\beta=\sum_iA_iD(\gamma^iP)\).
In the interpolation distribution, change coordinates by
\(\gamma^{-i}\) on the orbit of weight \(\kappa^i\). This is
well-defined since \(\kappa\bmod q\) and \(\gamma\) both have order
\(\ell^{\nu-c}\). The designated character block then has divisor
\(\operatorname{div}\beta\); the other total block divisors vanish.
Steinitz transport chooses matrices carrying the corresponding
lattices with determinant \(\beta\) on that block and \(1\) elsewhere.

Let \(\sum n_t(d_t)_*z=\partial b'\) be the given relation, choosing
a unit coefficient \(n_{t_0}\). Apply the two lattice actions
\(\sigma,\rho\) to \(b'\) and close their difference with the sum of
transport prisms on \((\sigma d_t)_*z\), weighted by \(n_t\). The
prism boundary relation cancels the difference term by term.
Divide the resulting cycle by \(n_{t_0}\) and add a trivial
complement to the projective lattice to place it in a free lattice.
The complement contributes zero in positive-degree homology.

Choose one denominator below every prime in the support of \(D\)
and clearing the kernel and cokernel of the integral map underlying
\(F\otimes_LF\simeq\prod_{\tau\in\operatorname{Gal}(F/L)}F\).
This makes all conjugates of the required elements units, splits the
coordinates for every allowed reduction, and equalizes the localized
lattices. If \(\omega(\tau)\) is the exponent of \(\tau\) on \(\zeta\),
the scalar evaluation and the determinant calculation now give
\(S(\beta)\), where \(S(t)=\sum_\tau\omega(\tau)^E v(r(\tau t))\).
Reindexing yields \(S(\gamma^{-i}\alpha)=a^iS(\alpha)\). Hence
\(S(\beta)=p(a)S(\alpha)=p(a)v(r(N_a(\alpha)))\), with the last
residue interpreted in the common localization. The construction of
\(y,b\) precedes \(r,v,\theta\).
\end{outline}
